% =====================================================================
\chapter{Problemas}
\label{chap:parteB}
% =====================================================================

La parte~B consta de seis ejercicios con un total de 80~puntos. La puntuación de cada apartado figura entre paréntesis. Debe justificarse cada paso y expresarse cada resultado con sus unidades.

% ---------------------------------------------------------------------
\section*{Módulo 14. El bloque decoder}

\begin{ejercicio}{2}{eje:B-bloque}
\textbf{(15 puntos.)} Un Transformer decoder-only del estilo de LLaMA recibe $n$ tokens y tiene dimensión de modelo $d$, $L$ bloques, $h$ cabezas de atención, MLP SwiGLU de dimensión oculta $\dff$, normalización RMSNorm, codificación de posición RoPE y un vocabulario de $V$ tokens con embeddings sin atar.
\begin{enumerate}[label=(\alph*)]
  \item \pts{5} Dibuje el diagrama completo, desde los tokens enteros hasta la distribución del token siguiente, con un bloque detallado. Señale las dos conexiones residuales y las dos normalizaciones del bloque, el lugar en que entra la información de posición y la normalización final.
  \item \pts{3} Escriba las ecuaciones del bloque $\ell$, con entrada $h^{\ell-1}$ y salida $h^{\ell}$, y la ecuación de los logits.
  \item \pts{3} Indique la forma de: la entrada al primer bloque; $Q$ tras separar las cabezas; la matriz de atención de una cabeza; la salida de la proyección $W_{\mathrm{gate}}$; la salida de un bloque, y los logits.
  \item \pts{2} ¿Qué subcapa mezcla información \emph{entre} tokens y cuál procesa cada token \emph{por separado}? ¿Por qué cada subcapa debe devolver un vector de $\R^{d}$?
  \item \pts{2} Dé dos razones por las que los modelos de lenguaje actuales prefieren la variante pre-norm a la post-norm.
\end{enumerate}
\fuente{adaptado de Stanford CS336~\cite{cs336}, tarea~1, §3.4 y §3.5, ecuación~(15) y problemas \emph{transformer\_block} y \emph{transformer\_lm}, y §7.3, problema \emph{pre\_norm\_ablation} (la implementación y el experimento se sustituyen por el diagrama, las ecuaciones y el análisis).}
\end{ejercicio}

% ---------------------------------------------------------------------
\section*{Módulo 13. Atención con máscara causal}

\begin{ejercicio}{2}{eje:B-atencion}
\textbf{(15 puntos.)} Tres tokens tienen embeddings de dimensión $d=2$: $x_1=(1,0)$, $x_2=(0,1)$ y $x_3=(1,1)$, apilados como filas de una matriz $X\in\R^{3\times 2}$. Una cabeza de atención usa $d_k=4$, $d_v=2$ y
\[
W_Q=\begin{pmatrix}1&1&0&0\\0&0&1&1\end{pmatrix},\qquad
W_K=\begin{pmatrix}2&0&0&0\\0&0&2&0\end{pmatrix},\qquad
W_V=\begin{pmatrix}1&1\\0&1\end{pmatrix}.
\]
\begin{enumerate}[label=(\alph*)]
  \item \pts{3} Calcule $Q=XW_Q$, $K=XW_K$ y $V=XW_V$.
  \item \pts{3} Calcule $S=QK\T$ y $S/\sqrt{d_k}$.
  \item \pts{2} Escriba la matriz de máscara causal $M$ y la matriz $Z=S/\sqrt{d_k}+M$. ¿Por qué la máscara usa $-\infty$ y no~$0$?
  \item \pts{3} Calcule $A=\softmax(Z)$, aplicada por filas, y la salida $O=AV$.
  \item \pts{2} Se sustituye $x_3$ por $(5,-3)$. Sin hacer cuentas largas, indique qué filas de $O$ cambian, justificándolo con la estructura de $A$, y calcule la nueva fila $o_3$. ¿Qué hace la cabeza en este caso?
  \item \pts{2} Calcule $\partial a_{22}/\partial z_{22}$ con los valores del apartado~(d). ¿Qué ocurriría con esta derivada si no se dividiera por $\sqrt{d_k}$ y las puntuaciones fueran muy grandes?
\end{enumerate}
\fuente{adaptado de Stanford CS224N~\cite{cs224n}, tarea~3, pregunta~1(a), «Copying in attention» (apartado~e), y de Jurafsky y Martin~\cite{jm3}, cap.~7, ejercicio~7.2 (apartado~c); datos numéricos propios.}
\end{ejercicio}

% ---------------------------------------------------------------------
\section*{Módulo 12. BPE a mano}

\begin{ejercicio}{2}{eje:B-bpe}
\textbf{(12 puntos.)} Un corpus contiene estas palabras con sus frecuencias: \tok{pino}~$\times 4$, \tok{pinos}~$\times 2$, \tok{vino}~$\times 3$, \tok{vinos}~$\times 1$ y \tok{pinta}~$\times 2$. Los símbolos iniciales son los caracteres y las fusiones no cruzan los límites entre palabras.
\begin{enumerate}[label=(\alph*)]
  \item \pts{1} ¿Cuántos símbolos tiene el corpus al empezar?
  \item \pts{6} Aplique cuatro fusiones BPE. En cada paso indique el par elegido y su frecuencia, el par que queda en segundo lugar y el número total de símbolos del corpus tras la fusión.
  \item \pts{3} Codifique con las cuatro reglas obtenidas las palabras \tok{vinos}, \tok{pinto}, \tok{sino} y \tok{ovino}, e indique el número de tokens de cada una.
  \item \pts{2} Un tokenizador BPE a nivel de byte procesa la palabra «piñón». ¿Cuántos símbolos iniciales tiene? ¿Qué pares de bytes conviene fusionar pronto en un corpus en español y por qué?
\end{enumerate}
\fuente{adaptado de Stanford CS336~\cite{cs336}, tarea~1, §2.4, ejemplo \emph{bpe\_example}, y §2.6.1, ejemplo \emph{bpe\_encoding}, y de Jurafsky y Martin~\cite{jm3}, §2.4.1 y ejercicio~2.3 (corpus y palabras cambiados).}
\end{ejercicio}

% ---------------------------------------------------------------------
\section*{Módulo 14. Parámetros, FLOPs y memoria}

\begin{ejercicio}{3}{eje:B-lince}
\textbf{(15 puntos.)} El modelo ficticio \textbf{Lince-1B} tiene $d=2048$, $L=24$, $h=16$, MLP SwiGLU, $V=48\,000$, embeddings \emph{atados}, RMSNorm con solo vector de pesos, RoPE y ningún sesgo.
\begin{enumerate}[label=(\alph*)]
  \item \pts{2} Calcule $\tfrac{8}{3}d$ y redondéelo al múltiplo de 256 inmediatamente superior para obtener $\dff$.
  \item \pts{4} Calcule el número de parámetros de un bloque y el total~$N$. Compare $N$ con la aproximación $12Ld^{2}$ e indique qué fracción de $N$ corresponde a los MLP.
  \item \pts{3} Estime los FLOPs por token en la generación y los FLOPs del entrenamiento con $D=2\cdot 10^{11}$ tokens. ¿Cuántas horas de GPU H100 ($989\cdot 10^{12}$~FLOP/s en bf16) se necesitan con un aprovechamiento del 40\,\%? ¿Cuántas horas de reloj con 64~GPU?
  \item \pts{3} Calcule la memoria que ocupan los pesos en bf16 y en int4, y la del estado del entrenamiento con Adam en precisión mixta. ¿Cabe el entrenamiento en una GPU de 24~GB?
  \item \pts{3} La atención añade unos $4Lnd$ FLOPs por token con un contexto de $n$ tokens. ¿Qué fracción de $2N$ supone con $n=4096$? ¿Para qué $n$ iguala a $2N$?
\end{enumerate}
\fuente{adaptado de Stanford CS336~\cite{cs336}, tarea~1, §3.5, problema \emph{transformer\_accounting}, apartados (a) a~(e), y §4.3, problema \emph{adamw\_accounting}, apartado~(d), y de Jurafsky y Martin~\cite{jm3}, cap.~7, ejercicio~7.1 (se sustituye GPT-2~XL por un modelo ficticio y la precisión fp32 por la mixta).}
\end{ejercicio}

% ---------------------------------------------------------------------
\section*{Módulos 13 y 14. Posiciones y coste de la atención}

\begin{ejercicio}{2}{eje:B-rope}
\textbf{(11 puntos.)}
\begin{enumerate}[label=(\alph*)]
  \item \pts{5} Se aplica RoPE en dimensión $d_k=2$ con $\theta=\pi/4$, $q=(1,0)$ y $k=(1,1)$. Use la identidad
  \[
  q\T R(\varphi)\,k=(q\cdot k)\cos\varphi-(q_1k_2-q_2k_1)\sin\varphi,\qquad \varphi=(n-m)\theta,
  \]
  para calcular la puntuación entre la consulta en la posición~$m$ y la clave en la posición~$n$ en los pares $(m,n)=(2,1)$, $(5,4)$, $(3,3)$, $(1,2)$ y $(0,4)$. ¿Qué se observa en los dos primeros?
  \item \pts{2} En Lince-1B (\cref{eje:B-lince}), ¿cuánto vale $d_k$? Si se duplica $h$ manteniendo $d$, ¿cuántos parámetros de atención se ganan?
  \item \pts{4} Con un contexto de $n=32\,768$ tokens, calcule la memoria en fp16 de la matriz de atención de una cabeza y la de todas las cabezas y capas de Lince-1B. ¿A partir de qué $n$ dominan los FLOPs cuadráticos de la atención, $4n^{2}d$ por capa, sobre los de las proyecciones, $8nd^{2}$ por capa? ¿Qué técnica evita almacenar $A$ completa?
\end{enumerate}
\fuente{adaptado de Stanford CS336~\cite{cs336}, tarea~1, §3.4.3, ecuaciones~(8) y~(9) (RoPE), y §3.5, problema \emph{transformer\_accounting}, apartado~(e) (se pide el cálculo a mano en dimensión~2 y la memoria de $A$).}
\end{ejercicio}

% ---------------------------------------------------------------------
\section*{Módulo 15. Entrenar y leer un mini-GPT}

\begin{ejercicio}{2}{eje:B-minigpt}
\textbf{(12 puntos.)} Un mini-GPT de caracteres tiene $V=80$, contexto $T=128$, $d=128$, $L=4$, $h=4$, MLP clásico de dimensión oculta $4d$ con GELU, LayerNorm con peso y sesgo, posiciones aprendidas, embeddings atados y ninguna capa lineal con sesgo. Se entrena con lotes de $B=32$ secuencias durante 5000~pasos sobre un corpus de $10^{6}$ caracteres. El \cref{tab:curvas} recoge sus pérdidas medias.

\begin{table}[H]
\centering\small
\setlength{\tabcolsep}{3.6pt}
\caption{Pérdidas del mini-GPT, en nats por carácter.}\label{tab:curvas}
\begin{tabular}{@{}lccccccccccc@{}}
\toprule
Paso & 0 & 500 & 1000 & 1500 & 2000 & 2500 & 3000 & 3500 & 4000 & 4500 & 5000 \\
\midrule
Entrenamiento & 4{,}39 & 2{,}12 & 1{,}66 & 1{,}47 & 1{,}38 & 1{,}30 & 1{,}22 & 1{,}15 & 1{,}08 & 1{,}03 & 0{,}98 \\
Validación & 4{,}39 & 2{,}14 & 1{,}70 & 1{,}52 & 1{,}46 & 1{,}44 & 1{,}45 & 1{,}48 & 1{,}52 & 1{,}57 & 1{,}62 \\
\bottomrule
\end{tabular}
\end{table}

\begin{enumerate}[label=(\alph*)]
  \item \pts{3} Cuente los parámetros del modelo.
  \item \pts{2} Explique por qué la pérdida inicial vale unos $4{,}39$ nats.
  \item \pts{3} ¿Qué punto de control debe conservarse? Calcule su perplejidad de validación y sus bits por carácter. Un modelo de bigramas alcanza una pérdida de validación de $2{,}45$ nats: ¿cuál es su perplejidad?
  \item \pts{2} ¿Cuántos tokens procesa el entrenamiento? ¿A cuántas épocas equivale? Relacione el resultado con el \cref{tab:curvas} y proponga dos remedios.
  \item \pts{2} Al generar, el modelo produce los logits $(2,1,0)$ para tres caracteres. Calcule las probabilidades con temperatura $\tau=1$ y $\tau=0{,}5$. ¿Qué efecto tiene bajar la temperatura sobre las muestras?
\end{enumerate}
\fuente{adaptado de Stanford CS336~\cite{cs336}, tarea~1, §4.1, ecuación~(18) (perplejidad), y §6, ecuación~(23) y problema \emph{decoding} (temperatura), y de Jurafsky y Martin~\cite{jm3}, cap.~7, ejercicios~7.4 y~7.5 (curvas y logits propios).}
\end{ejercicio}
